\section{Discrepancy certificates from continued fractions (Ostrowski/Denjoy--Koksma)}
\label{app:discrepancy_ostrowski}
\noindent
This appendix makes explicit the quantitative link used in
Section~\ref{subsubsec:discrepancy_certificates}:
continued-fraction data of the scan slope $\alpha$ yield deterministic
finite-$N$ discrepancy bounds for the Kronecker orbit
$x_n=x_0+n\alpha\ (\mathrm{mod}\ 1)$.
We keep the discussion in the one-dimensional setting relevant to window
counts, where the key estimate reduces to interval indicators.
\par
\AuditTag Unification note: the same continued-fraction denominators $(q_k)$ also control \emph{return/persistence} scales for the rotation-window (Sturmian) coding.
See Appendix~\ref{app:sturmian_return_times_qk} for a two-scale return-time bound $\tau\le q_k+q_{k+1}$ for length-$m$ factors when $q_k\le m<q_{k+1}$, and Appendix~\ref{app:rotation_hitting_time_discrepancy} for a discrepancy-derived hitting-time bound at convergent scales.

\subsection{Star discrepancy and a rational baseline}
\label{app:discrepancy_def}
For a finite point set $P_N=\{x_0,\dots,x_{N-1}\}\subset[0,1)$, the
one-dimensional star discrepancy is
\[
D_N^\ast(P_N):=\sup_{a\in[0,1]}
\left|
\frac{1}{N}\#\{0\le n\le N-1:\ x_n<a\}-a
\right|.
\]

\begin{lemma}[Equally spaced points have discrepancy $1/N$]
\label{lem:equally_spaced_discrepancy}
Let $y_n=\{y_0+n/N\}$ for $n=0,\ldots,N-1$.
Then $D_N^\ast(\{y_0,\ldots,y_{N-1}\})=1/N$.
\end{lemma}
\begin{proof}
The set $\{y_n\}$ is a translate of the uniform grid $\{j/N\}_{j=0}^{N-1}$,
and star discrepancy is translation invariant on the circle.
For the grid, for any $a\in[0,1]$ the count
$\#\{j/N<a\}$ equals either $\lfloor Na\rfloor$ or $\lceil Na\rceil-1$,
so the normalized error differs from $a$ by at most $1/N$.
Taking $a=j/N$ shows the bound is attained, hence $D_N^\ast=1/N$.
\end{proof}

\subsection{Continued fractions and convergents}
\label{app:continued_fractions}
Let $\alpha\in(0,1)\setminus\QQ$ with continued fraction
$\alpha=[0;a_1,a_2,\ldots]$ and convergents $p_k/q_k$.
We use the standard recurrences
\[
q_{-1}=0,\ q_0=1,\qquad q_{k+1}=a_{k+1}q_k+q_{k-1}\quad (k\ge 0),
\]
and similarly for $p_k$.

\begin{lemma}[Best-approximation error of convergents]
\label{lem:convergent_error}
For every $k\ge 0$,
\[
\left|\alpha-\frac{p_k}{q_k}\right|
<\frac{1}{q_k q_{k+1}}
\le \frac{1}{q_k^2}.
\]
\end{lemma}
\begin{proof}
This is a standard continued-fraction property; see, e.g.,
\cite[Ch.\ 10]{HardyWright2008} or \cite{KuipersNiederreiter1974}.
\end{proof}

\subsection{A discrepancy bound at convergent lengths}
\label{app:discrepancy_convergent_lengths}
Let $x_n=\{x_0+n\alpha\}$ be the scan orbit.
At convergent lengths $N=q_k$, the orbit is uniformly close to the rational
orbit with step $p_k/q_k$.
This yields a simple $O(1/N)$ discrepancy bound with an explicit constant.

\begin{lemma}[Convergent-length discrepancy bound]
\label{lem:convergent_length_discrepancy}
Fix $k\ge 0$ and set $N:=q_k$.
Let $P_N(\alpha)=\{x_0,\dots,x_{N-1}\}$ with $x_n=\{x_0+n\alpha\}$.
Then
\[
D_N^\ast(P_N(\alpha))\le \frac{8}{N}.
\]
\end{lemma}
\begin{proof}
Let $\beta:=p_k/q_k$ and define the rational reference points
$y_n:=\{x_0+n\beta\}$.
Since $\gcd(p_k,q_k)=1$, the set $\{y_n\}_{n=0}^{N-1}$ is exactly an equally
spaced grid (a permutation of $\{x_0+j/N\}$), so by
Lemma~\ref{lem:equally_spaced_discrepancy} one has $D_N^\ast(\{y_n\})=1/N$.

By Lemma~\ref{lem:convergent_error}, $\delta:=|\alpha-\beta|<1/N^2$.
For each $n\in\{0,\ldots,N-1\}$, the points $x_n$ and $y_n$ differ by the
rotation $n(\alpha-\beta)$ on the circle, so their \emph{circular distance}
obeys
\[
d_{\TT}(x_n,y_n)\le n\,\delta < \frac{1}{N},
\qquad
d_{\TT}(u,v):=\min\{|u-v|,\,1-|u-v|\}.
\]

Fix $a\in[0,1]$ and compare the counts
$C_x(a):=\#\{0\le n\le N-1:\ x_n<a\}$ and
$C_y(a):=\#\{0\le n\le N-1:\ y_n<a\}$.
If $x_n$ and $y_n$ fall on different sides of the interval boundary
of $[0,a)$, then $y_n$ must lie within circular distance $<1/N$ of one of the
two boundary points $\{0,a\}$ (otherwise a $<1/N$ perturbation cannot change
membership in the half-open interval).
Since the $y_n$ are $1/N$-spaced, there are at most $3$ indices $n$ within
circular distance $<1/N$ of each boundary point, hence
$|C_x(a)-C_y(a)|\le 6$ for every $a$.

Hence
\[
\left|\frac{C_x(a)}{N}-a\right|
\le
\left|\frac{C_y(a)}{N}-a\right|
\;+\;
\frac{|C_x(a)-C_y(a)|}{N}
\le \frac{1}{N}+\frac{6}{N}=\frac{7}{N},
\]
uniformly over $a$.
Taking the supremum gives $D_N^\ast(P_N(\alpha))\le 7/N$.
We state the slightly looser constant $8/N$ for a clean margin.
\end{proof}

\subsection{Ostrowski decomposition and a general finite-\texorpdfstring{$N$}{N} bound}
\label{app:ostrowski_discrepancy}
The convergent-length bound uplifts to arbitrary $N$ using the Ostrowski
representation, which decomposes $N$ into a sum of convergent denominators.

\begin{definition}[Ostrowski representation (statement)]
\label{def:ostrowski_representation_statement}
Let $\alpha=[0;a_1,a_2,\ldots]$ be irrational with convergent denominators
$q_k$.
Every integer $N\ge 1$ admits a (unique) representation
\[
N=\sum_{k=0}^{m} b_k q_k,
\qquad
0\le b_k\le a_{k+1},
\]
with the usual local admissibility constraint
$b_k=a_{k+1}\Rightarrow b_{k-1}=0$ for $k\ge 1$.
\end{definition}

\begin{proposition}[Digit-sum discrepancy bound]
\label{prop:digit_sum_discrepancy}
Let $x_n=\{x_0+n\alpha\}$ and let $P_N(\alpha)=\{x_0,\ldots,x_{N-1}\}$.
Write $N=\sum_{k=0}^{m} b_k q_k$ in Ostrowski form.
Then
\[
D_N^\ast(P_N(\alpha))
\le
\frac{8}{N}\sum_{k=0}^{m} b_k.
\]
\end{proposition}
\begin{proof}
Decompose the prefix $\{0,1,\ldots,N-1\}$ into consecutive blocks of lengths
$q_k$, repeated $b_k$ times for each $k$ (largest $k$ to smallest).
For each such block of length $q_k$, apply
Lemma~\ref{lem:convergent_length_discrepancy} with the intercept shifted to the
block start (i.e.\ replace $x_0$ by $x_t$ for the corresponding start index
$t$).
This bounds the star discrepancy of each length-$q_k$ block by $8/q_k$,
equivalently its \emph{count error} (unnormalized) by at most $8$.
Summing count errors over all blocks yields a total count error bounded by
$8\sum_k b_k$ uniformly in the threshold $a$.
Dividing by $N$ and taking the supremum in the definition of $D_N^\ast$
gives the stated bound.
\end{proof}

\subsection{Bounded type and why the golden branch is minimax}
\label{app:bounded_type}

\begin{remark}[Standard bounded-type discrepancy bound (reference)]
\label{rem:bounded_type_standard_reference}
For an irrational rotation $x_n=\{x_0+n\alpha\}$, if the continued-fraction digits of $\alpha$ are bounded, $a_k\le A$ (``bounded type''), then the Kronecker sequence has star discrepancy of order $\log N/N$ with a constant depending only on $A$; see \cite{KuipersNiederreiter1974,Cassels1957}.
The explicit digit-sum certificate recorded below is a self-contained specialization to the one-dimensional interval-indicator case with an auditable constant and a transparent dependence on $A$.
\end{remark}
\begin{corollary}[Bounded-type logarithmic rate]
\label{cor:bounded_type_discrepancy}
Assume $\alpha$ has bounded partial quotients: $a_k\le A$ for all $k$.
Then for every $N\ge 1$,
\[
D_N^\ast(P_N(\alpha))
\le
\frac{8A}{N}\,\Bigl(4+\log_\varphi N\Bigr),
\qquad
E_N:=N D_N^\ast(P_N(\alpha))\ \le\ 8A\,(4+\log_\varphi N).
\]
\end{corollary}
\begin{proof}
Let $N=\sum_{k=0}^m b_k q_k$ be the Ostrowski representation.
Since $b_k\le a_{k+1}\le A$, one has $\sum_{k=0}^m b_k\le A(m+1)$.

Next, since $a_{k+1}\ge 1$ for all $k$, the denominators satisfy
$q_{k+1}=a_{k+1}q_k+q_{k-1}\ge q_k+q_{k-1}$.
With $q_0=1$ and $q_1\ge 1$, this implies $q_m\ge F_{m+1}$ for all $m\ge 1$.
By Binet's formula, $F_{m+1}\ge \varphi^{m-1}/\sqrt{5}$ for $m\ge 1$, hence
if $q_m\le N$ then
\[
m-1\le \log_\varphi(\sqrt{5}\,N)\le 2+\log_\varphi N.
\]
Therefore $m+1\le 4+\log_\varphi N$.

Combining with Proposition~\ref{prop:digit_sum_discrepancy} yields
\[
D_N^\ast(P_N(\alpha))
\le \frac{8}{N}\sum_{k=0}^{m} b_k
\le \frac{8A(m+1)}{N}
\le \frac{8A}{N}\bigl(4+\log_\varphi N\bigr),
\]
and multiplying by $N$ gives the bound on $E_N$.
\end{proof}

\begin{remark}[Golden branch as a minimax bounded-type choice]
\label{rem:golden_minimax_discrepancy}
The constant in Corollary~\ref{cor:bounded_type_discrepancy} depends
monotonically on the bound $A$ on continued-fraction digits.
The golden branch is characterized by $a_k\equiv 1$, i.e.\ it is the unique
irrational of constant type with the minimal possible bound $A=1$.
Thus, within the bounded-type class, it is the canonical minimax choice for
discrepancy certificates at finite $N$:
it gives the smallest explicit worst-case upper bound in the audited family
of certificates derived from digit sums.
\end{remark}


